\uuid{PlrU}
\titre{Strucutre de sous-espace vectoriel}
\auteur{hanani ; mzari}
\organisation{exo7}
\licence{CC BY-NC-SA 4.0}
\datecreate{2020-01-20}
\niveau{L1}
\module{Algèbre}
\chapitre{Espace vectoriel}
\sousChapitre{Définition, sous-espace}
\difficulte{2}
\theme{Espaces vectoriels, Strucutre de sous-espace vectoriel}
\exo7qcmid{lille-2.espaces-vectoriels.005}

\contenu{
\begin{question}
\qclassification{106.01}
\qtags{motcle=Strucutre de sous-espace vectoriel}
Soit $E=\{(x,y,z) \in \Rr^3 \, ; \;  xy+xz+yz = 0\}$, muni des opérations usuelles. Quelles sont les assertions vraies ?
\begin{answers}
\good{$(0,0,0)\in E$.}
\good{$E$ n'est pas stable par addition.}
\good{$E$ est stable par multiplication par un scalaire.}
\bad{$E$ est un sous-espace vectoriel de $\Rr^3$.}
\end{answers}
\begin{explanations}
$E$ n'est pas un sous-espace vectoriel de $\Rr^3$, puisque $E$ n'est pas stable par addition :
$(1,0,0), (0,1,0) \in E$, mais $(1,1,0) \notin E$. Cependant, $E$ est stable par multipliction par un scalaire.
\end{explanations}
\end{question}
}
